\begin{definition}
    Преобразование $\phi$ называется диагонализируемым, если существует базис, в котором его матрица имеет диагональный вид.
\end{definition}
\begin{theorem}
    Преобразование $A$ диагонализируемо тогда и только тогда, когда все пространство $L=L_1\oplus\dots\oplus L_k$ -- прямая сумма собственных подпространств.
\end{theorem}
\begin{proof} 
$\Leftarrow$ Если $L=L_1\oplus\dots\oplus L_k$, то объединение базисов из $L_i$ будет базисом в $L$. Тогда в этом базисе по теореме 7.1 матрица $A$ -- диагональная.

$\Rightarrow$ Если матрица $A$ -- диагональная, то по теореме 7.1 существует базис из собственных векторов. Сгруппируем эти векторы по собственным числам $\lambda_i$, которым они соответствуют. Тогда каждому из этих собственных чисел соответствует собственное подпространство $L_i$. А так как  объединение базисов из $L_i$ -- базис в $L$, то сумма этих собственных подпространств прямая и $L_1\oplus\dots\oplus L_k=L$.
\end{proof}
Приведем необходимые и достаточные условия диагонализируемости $A$:
\begin{proposition} [Необходимое условие]
    Если преобразование $A$ диагонализируемо, тогда все корни характеристического многочлена $P(\lambda)$ вещественные.
\end{proposition}
\begin{proof}
    Пусть $\lambda_i\in \R$ корень характеристического многочлена с алгебраической кратностью  $S_i$, тогда по теореме 6.4 $\dim L_i\leqslant S_i$, где $L_i$ соответсвующее собственное подпространство. Тогда $\sum\dim L_i \leqslant \sum S_i\leqslant \deg P(\lambda)\leqslant n$. Но если есть корень $\lambda^{*}\in \mathbb{C}$, то $\sum S_i< \deg P(\lambda)$. Следовательно $\sum\dim L_i <n$, а значит не существует базиса из собственных векторов, что противоречит теореме 7.1.
\end{proof}
\textit{Замечание:} недостаточность данного условия демонстрирует пример преобразования с матрицей $A=\left(\begin{matrix}
    1&1\\0&1
\end{matrix}\right)$, так как в данном случае алгебраическая кратность корня не равна геометрической кратности.
\begin{proposition} [Достаточное условие]
    Если все корни характеристического многочлена $P(\lambda)$ вещественные и различные, то преобразование $A$ диагонализируемо.
\end{proposition}
\begin{proof}
    Для каждого корня $\lambda_i$ выберем вектор $f_i$ из соответствующего собственного подпространства. По утверждению 6.1 набор $f_1, \dots, f_n$ ЛНЗ. Тогда его можно взять в качестве базиса, а по теореме 7.1 в нем матрица $A$ диагональная.
\end{proof}
Это достаточное условие можно переформулировать в матричном виде:
\begin{theorem}
    Если все собственные значения преобразования $A$ вещественные и различные, то существует невырожденная матрица $S$ такая, что $S^{-1}AS$ диагональная.
\end{theorem}
Усилив утверждение 7.1 можно получить критерий диагонализируемости:
\begin{theorem} [Критерий диагонализируемости]
    Преобразование диагонализируемо тогда и только тогда, когда все корни $\lambda_1, \dots , \lambda_n$ характеристического многочлена $P(\lambda)$ вещественные и для всех $\lambda_i$ алгебраическая кратность равна геометрической.
\end{theorem}
\begin{proof}
    TBA...
\end{proof}
\begin{theorem} [Гамильтона-Кэли]
    Пусть $P(\lambda)$ -- характеристический многочлен матрицы $A$. Тогда $P(A)=0$.
\end{theorem}
\begin{proof}
    Выберем $\lambda$ не являющееся корнем $P(\lambda)$. Тогда матрица $A-\lambda E$ невырожденная, а тогда существует $$(A-\lambda E )^{-1}=\frac{1}{\det (A-\lambda E)}\begin{pmatrix}
        (-1)^{i+j}D_{ji}
    \end{pmatrix}=\frac{1}{\det (A-\lambda E)}\begin{pmatrix}
        B(\lambda)
    \end{pmatrix} \ \ \ \ (*)$$
    где $B(\lambda)$ -- элементы-многочлены степени не выше $n-1$. Значит $B(\lambda)=B_0+\dots+B_{n-1}\lambda^{n-1}$. Домножим уравнение $(*)$ на $(A-\lambda E)\det(A-\lambda E)$:
    \[
    P(\lambda)E=(A-\lambda E)^{-1}(A-\lambda E)\det(A-\lambda E)=(A-\lambda E)B(\lambda)
    \]
    Тогда расписав $P(\lambda)=a_n \lambda^nE+\dots+a_0E$ получаем:
    \[
    a_n \lambda^nE+\dots+a_0E=(A-\lambda E)(B_0+\dots+B_{n-1}\lambda^{n-1})
    \]
    Это равенство верно при всех $\lambda$, не являющихся корнями $P(\lambda)$. Следовательно, эти два многочлена совпадают, а значит имеем условие на равенство коэффициентов:
    \begin{align*}
    a_0 E &= AB_0  \ \ \ &|\cdot A^0\\
    a_1 E &= AB_1 - B_0 E\ \ \ &|\cdot A^1\\
    a_2 E &= AB_2 - B_1 E \ \ \ &|\cdot A^2\\
    \vdots \\
    a_n E&=-B_{n-1} \ \ \ &|\cdot A^n
    \end{align*}
    Домножая каждое равенство на соответствующую степень матрицы $A$ и складывая все равенства справа получаем:
    \[
    a_n A^n+\dots+a_1A+a_0E=0=P(A)
    \]
\end{proof}

\section{Линейные функции}
\textit{(данная тема не будет присутствовать в билетах, но спросить определения могут)}
\begin{definition}
    Функция $f:L\to \R$, заданная на линейном пространстве $L$, называется линейной, если
    \begin{align*}
    &1) \ \forall x, y \in L \hookrightarrow f(x+y)=f(x)+f(y) \\
    &2)\  \forall \alpha \in \R \ \forall x\in L \hookrightarrow f(\alpha x)=\alpha f(x)
    \end{align*}
\end{definition}
То есть линейная функция -- частный случай линейного отображения и все утверждения про линейные отображения, сформулированные ранее, остаются верными и для линейных функций.

Пусть в $L$ задан базис $e^{(n)}$. Тогда 
\[
f(x)=f(x_1e_1+\dots +x_ne_n)=x_1f(e_1)+\dots +x_nf(e_n)=(f(e_1), \dots, f(e_n))\begin{pmatrix} x_1 \\ \vdots \\ x_n \end{pmatrix}
\]
где сомножитель $(f(e_1), \dots, f(e_n))$ называется координатной строкой функции.

При переходе к базису $e\to e'$ (посредствам матрицы $S$) координатная строка $F\to F'S$.

Аналогично отображениям на множестве линейных функций $L\to \R$ вводим сумму функций и умножение функций на скаляр:
\begin{align*}
    &\forall f,g \ \forall x\in L \hookrightarrow(f+g)(x)=f(x)+g(x) \\
    &\forall f \ \forall \alpha\in \R \ \forall x\in L \hookrightarrow (\alpha f)(x)=\alpha f(x)
\end{align*}
Тогда множество всех линейных функций из $L$, с введенными операциями является линейным пространством, \textit{сопряженным} к $L$, обозначается $L^*$.

В качестве базиса в $L^*$ выберем функции $f_i$ такие, что $f_i(x)=x_i$ ($i$-ая координата вектора $x\in L$). Тогда $$f_i(e_j)=\begin{cases} 
   1, \ \ i=j \\
   0, \ \ i \neq j 
\end{cases}$$ Такой базис называется биортогональным (взаимным) к базису $e$.
Также следствием построения такого базиса является $\dim L^*=\dim L$.

\section{Билинейные и квадратичные формы}
\subsection{Билинейные формы}
\begin{definition}
    Функция $b(x,y):L\times L\to \R$ называется билинейной, если она линейна по каждому аргументу:
    \begin{align*}
        &\forall x, y, z \in L \hookrightarrow b(x+y,z)=b(x,z)+b(y,z), \ \ b(x, y+z)=b(x,y)+b(x,z)\\
        &\forall x,y\in L \ \forall \alpha \in \R \hookrightarrow b(\alpha x, y)=\alpha b(x,y), \ \ b(x, \alpha y)=\alpha b(x,y)
    \end{align*}
\end{definition}
\textit{Замечание:} билинейная форма $\equiv$ билинейная функция.

Пусть в $L$ задан базис $e^{(n)}$. Тогда 
\begin{align*}
    x &=x_1e_1+\dots +x_ne_n \\
    y &= y_1e_1+\dots+y_ne_n 
\end{align*}
\[
b(x,y)=b(x_1e_1+\dots +x_ne_n, y_1e_1+\dots+y_ne_n)=\sum_{i,j=1}^nx_iy_jb(e_i,e_j)=(x_1, \dots, x_n)(B)\begin{pmatrix}
    y_1 \\
    \vdots \\
    y_n
\end{pmatrix}
\]
$$B=\begin{pmatrix}
    b(e_1,e_1) & \dots & b(e_1,e_n) \\
    \vdots & \ddots & \vdots \\
    b(e_n,e_1) & \dots & b(e_n, e_n)
\end{pmatrix}$$

Матрица $B$ называется \textit{матрицей билинейной формы}. Значит $b(x,y)=x^TBy$.
